\section{Интеграл Римана}

\subsection{Введение}

\begin{defn}{}{}
    Пусть задан отрезок \( [a, b] \)

    \begin{itemize}
        \item
            \(a = x_0 < x_1 < \ldots < x_{n - 1} < x_n = b \)

            И \( T = \{ x_0, \ldots, x_n \} \) называется разбиением \( [a, b] \)
        \item
            \( \Delta_k = [x_{k - 1}, x_k] \) называется \( k \)-м отрезком разбиения.

            \( |\Delta_k| = x_k - x_{k - 1} \) --- его длина.
        \item
            \( d(t) = \max\limits_{k} |\Delta_k| \) называется диаметром разбиения \( T \).
        \item
            \( \xi = \{ \xi_k : \xi_k \in \Delta_k \} \) называется разметкой разбиения \( T \)
            (взяли из каждого отрезка разбиения какую-то точку).
        \item
            Пусть на \( [a, b] \) задана \( f : [a, b] \to \RR \).
            Тогда \( \sigma(f, T, \xi) = \sum\limits_{k = 1}^{n} f(\xi_k) |\Delta_k| \)
            называется интегральной суммой или суммой Римана,
            соответственно разметки разбиения \( (T, \xi) \) и функции \( f \).
    \end{itemize}
\end{defn}

\begin{defn}{Интеграл Римана}{}
    Если
    \(
        \exists I \in \RR : \forall \varepsilon > 0 \ \exists \delta > 0 :
        \forall T : d(T) < \delta \ \forall \xi \
        |\sigma(f, T, \xi) - I| < \varepsilon
    \),
    то \( I \) называется интегралом функции \( f \) на \( [a, b] \):
    \[
        I = \int\limits_{a}^{b} f(x) dx
    \]
\end{defn}

\begin{example}{}{}
    \[
        f(x) =_{[a, b]} c
    \]
\end{example}
\[
    \sigma(f, T, \xi)
    =
    \sum\limits_{k = 1}^{n} C|\Delta_k|
    =
    C(b - a)
    =
    \int\limits_{a}^{b} f(x) dx
\]

\begin{exercise}{}{}
    \( f \) на \( [a, b] \) разбивается на конечное число отрезков,
    на которых она константа. Найдите ее интеграл.
\end{exercise}

\begin{example}{}{}
    \[
        f(x) = x^2, \quad x \in [0, 1]
    \]
\end{example}

Пусть все отрезки разбиения равны друг другу,
то есть \( \Delta_k = \frac{1}{n} \).
Пусть \( \xi_k = x_k \):
\[
    \sigma(f, T, \xi)
    =
    \sum\limits_{k = 1}^{n} \left( \frac{k}{n} \right)^2 \cdot \frac{1}{n}
    =
    \frac{1^2 + 2^2 + \ldots + n^2}{n^3}
    =
    \frac{n(n + 1)(2n + 1)}{n^3}
    \xmapsto{n \to \infty}
    \frac{1}{3}
\]

\begin{thrm}{Необходимое условие интегрируемости}{}
    \( f \in R[a, b] \Rightarrow f \) ограничена на \( [a, b] \)
\end{thrm}

Пусть \( f \) неограничена на \( [a, b] \).
Пусть \( T \) --- некоторое разбиение.
Пусть \( f \) неограничена \( \Delta_r \).
При выбранной разметке \( \xi \):
\[
    \sigma(f, T, \xi)
    =
    \sum\limits_{\substack{k = 1 \\ k \neq r}}^{n} f(\xi_k) |\delta_k| + f(\xi_r) |\xi_r|
\]

Зафиксируем все отрезки \( \Delta_k \) и все точки разметки кроме \( \xi_r \).

Пусть \( \sum\limits_{\substack{k = 1 \\ k \neq r}}^{n} f(\xi_k) |\delta_k| = A \).
Докажем, что выбором \( \xi_r \) \( |\sigma(f, T, \xi)| \) может быть сделан сколько угодно большим:
\[
    |\sigma(f, T, \xi)| \geq |f(\xi_r) | \delta_r | | - |A|
\]

Так как \( f \) неограничена на \( \delta_r \),
можно сделать \( f(\xi_r) \) сколь угодно большим.
А значит можно получить сколько угодно большую интегральную сумму.

\begin{example}{}{}
    \[
        D(x) = \begin{cases}
            1, x \in \QQ \cap [0, 1]
            \\
            0, x \in \II \cap [0, 1]
        \end{cases}
    \]

    Интеграл взять нельзя (берем разметки только из рациональных и только из иррациональных чисел).
\end{example}

\subsection{Критерий Дарбу}

\begin{defn}{Суммы Дарбу}{}
    Пусть \( f \) задана и ограничена на \( [a, b] \), \( T \) --- разбиение.

    \begin{itemize}
        \item
            \( s(T) = \sum\limits_{k = 1}^{n} m_k |\Delta_k| \),
            где \( m_k = \inf\limits_{x \in \Delta_k} f(x) \) называют нижней суммой Дарбу.
        \item
            \( S(T) = \sum\limits_{k = 1}^{n} M_k |\Delta_k| \),
            где \( M_k = \sup\limits_{x \in \Delta_k} f(x) \) называют верхней суммой Дарбу.
    \end{itemize}
\end{defn}

\begin{defn}{Интегралы Дарбу}{}
    \begin{itemize}
        \item
            \( \sup\limits_{T} s(T) = \ul{I} \) называют нижним интегралом Дарбу
        \item
            \( \inf\limits_{T} S(T) = \ol{I} \) называют верхним интегралом Дарбу
    \end{itemize}
\end{defn}

\begin{thrm}{Свойства сумм Дарбу}{}
    Пусть \( f \) ограничена на \( [a, b] \). Тогда
    \begin{enumerate}
        \item
            \( \forall T \ \forall \xi \quad s(T) \leq \sigma(f, T, \xi) \leq S(T) \)
        \item
            Если задано разбиение \( T \), то
            \begin{itemize}
                \item
                    \( \inf\limits_{\xi} \sigma(f, T, \xi) = s(T) \)
                \item
                    \( \sup\limits_{\xi} \sigma(f, T, \xi) = S(T) \)
            \end{itemize}
        \item
            \( \forall T_1, T_2 \quad s(T_1) \leq S(T_2) \)
        \item
            \( \ul{I}, \ol{I} \) существуют, причем \( s(T) \leq \ul{I} \leq \ol{I} \leq S(T) \)
        \item
            \( S(T) - s(T) \geq \ol{I} - \ul{I} \geq 0 \)
    \end{enumerate}
\end{thrm}

\begin{enumerate}
    \item
        \(
            \forall \xi_k \in \Delta_k
            \quad
            m_k | \Delta_k | \leq f(\xi_k) | \Delta_k | \leq M_k | \Delta_k |
        \)

        Просуммируем по всем \( k \), получим искомое.
    \item
        Так как \( m_k = \inf\limits_{\Delta_k} \), то
        \( \forall \varepsilon > 0 \ \exists \xi_k \in \Delta_k : f(\xi_k) < m_k + \varepsilon \)

        \(
            s(T) = \sum\limits_{k = 1}^{n} f(\xi_k) | \Delta_k |
            \leq
            \sum\limits_{k = 1}^{n} (m_k + \varepsilon) | \Delta_k |
            =
            s(T) + \varepsilon (b - a)
        \)

        Значит \( s(T) = \inf\limits_{\xi} \sigma(f, T, \xi) \) по определению, с \( S(T) \) так же.
    \item
        Если к разбиению добавить точки, то \( s(T) \uparrow \), а \( S(T) \downarrow \).

        Если \( T_3 = T_1 \cup T_2 \), то \( s(T_1) \leq s(T_3) \leq S(T_3) \leq S(T_2) \).
    \item
        \( s(T) \) ограничено сверху любой \( S(T') \),
        \( \Rightarrow \exists \ul{I} = \sup\limits_T s(T) \).
        Аналогично с \( \ol{I} \).

        Очевидно, что \( s(T) \leq \ul{I} \leq \ol{I} \leq S(T) \).
    \item
        Явно следует из предыдущего пункта.
\end{enumerate}

\begin{thrm}{Критерий Дарбу}{}
    \(
        f \in R[a, b]
        \Leftrightarrow
        \forall \varepsilon > 0 \ \exists \delta > 0 : \forall T \ d(T) < \delta \ S(T) - s(T) < \varepsilon
    \)

    Иначе говоря, необходимо и достаточно \( \lim\limits_{d(T) \to 0} S(T) - s(T) = 0 \)
\end{thrm}

\begin{itemize}
    \item
        \( \Rightarrow \)
        \begin{gather*}
            \exists I \in \RR : \forall \varepsilon > 0 \ \exists \delta > 0 : \forall T : d(T) < \delta
            \forall \xi I - \frac{\varepsilon}{3} < \sigma(f, T, \xi) < I + \frac{\varepsilon}{3}
            \\
            \Downarrow
            \\
            I - \frac{\varepsilon}{e} \leq s(T) \leq S(T) \leq I + \frac{\varepsilon}{3}
            \\
            \Downarrow
            \\
            S(T) - s(T) < \varepsilon
        \end{gather*}
    \item
        \( \Leftarrow \)

        Если \( \lim\limits_{d(T) \to 0} S(T) - s(T) = 0 \), то \( \ul{I} = \ol{I} = I \) и \( s(T) \leq I \leq S(T) \).

        \(
            I, \sigma(f, T, \xi) \in [s(T), S(T)]
            \Rightarrow
            \forall \varepsilon \ \exists \delta > 0 : \forall T : d(T) < \delta \ \forall \xi \ |\sigma(f, T, \xi) - I| < \varepsilon
        \)
\end{itemize}

\begin{thrm}{Критерий в терминах колебаний}{}
    \( S(T) - s(T) = \sum\limits_{k = 1}^{n} (M_k - m_k) | \Delta_k | = \Omega(T) \) --- колебание

    \( f \in R[a, b] \Leftrightarrow \lim\limits_{d(T) \to 0} \Omega(T) \)
\end{thrm}

\section{Какие функции интегрируемы?}

\begin{thrm}{Эквивалентные условия интегрируемости}{}
    \( f \in R[a, b] \), если выполнено одно из трех эквивалентных условий:
    \begin{enumerate}
        \item
            \( \lim\limits_{d(T) \to 0} \Omega(T) = 0 \)
        \item
            \( \ul{I} = \ol{I} \)
        \item
            \( \inf\limits_T \Omega(T) = 0 \) --- инфимум-критерий
    \end{enumerate}
\end{thrm}

\begin{exercise}{}{}
    Докажи \( (3) \to (1) \)
\end{exercise}

\begin{exercise}{}{}
    Пусть \( T_n \) равномерное разбиение отрезка:
    \[
        x_0 = a, x_1 = a + \frac{b - a}{n}, x_2 = a + 2 \cdot \frac{b - a}{n}, \ldots, x_n = b
    \]

    Докажите, что \( f \in R[a, b] \Leftrightarrow \lim\limits_{n \to \infty} \Omega(T_n) = 0 \)
\end{exercise}

\begin{defn}{Колебание}{}
    Колебанием функции \( f \) на множестве \( A \) называется
    \[
        \omega(f, A) = \sup\limits_{x, y \in A} (f(x) - f(y))
    \]

    Колебанием \( f \) в точке \( x \) называется
    \[
        \omega(x) = \inf\limits_{\delta > 0} \omega (f, U_{\delta} (x)) = \lim\limits_{\delta \to 0} \omega(f, U_{\delta})
    \]

\end{defn}

\begin{exercise}{}{}
    \( f \in C(x) \Leftrightarrow \omega(x) = 0 \)
\end{exercise}

\begin{thrm}{}{}
    Пусть \( L_{\alpha} = \{ x : \omega(x) \geq \alpha \} \)

    Тогда \( L_{\alpha} \) --- замкнутое множество.
\end{thrm}

Пусть \( y \) --- предельная точка \( L_{\alpha} \)

\(
    \exists \{ y_n \} \subset L_{\alpha} : \lim\limits_{n \to \infty} y_n = y
    \Rightarrow
    \forall \delta > 0 \ \exists n \in \NN : y_n \in U_{\delta} (y)
\)

\(
    \exists \delta_1 > 0 : U_{\delta_1} (y_n) \subset U_{\delta} (y)
    \Rightarrow
    \omega(f, U_{\delta}(y)) \geq \omega(f, U_{\delta_1} (y_n)) \geq \alpha
\)

Отсюда следует, что
\(
    \inf\limits_{\delta > 0} \omega(f, U_{\delta}(y)) \geq \alpha
    \Rightarrow
    y \in L_{\alpha}
\)

\begin{defn}{}{}
    Множество \( L \) имеет меру Лебега нуль,
    если \( \forall \varepsilon > 0 \)
    существует не более чем счетное семейство интервалов \( \{ I_n \} \)
    таких, что \( L \subseteq \bigcup\limits_{n = 1}^{\infty} I_n \)
    и \( \sum\limits_{n = 1}^{\infty} | I_n | < \varepsilon \)

    Обозначется \( \mu(L) = 0 \)
\end{defn}

\begin{example}{}{}
    \begin{enumerate}
        \item
            \( L = \{ a \} \quad \mu(L) = 0 \)
        \item
            \( L = \{ a_1, a_2, \ldots, a_k \} \quad \mu(L) = 0 \)
        \item
            \( L = \{ a_k \}_{k = 1}^{\infty} \)

            Пусть \( \varepsilon > 0 \) задано, \( a_k \) покрыто интервалом \( I_k \),
            таким что \( | I_k | < \frac{\varepsilon}{2^{k + 1}} \).

            Тогда \( \sum\limits_{k = 1} |I_k| < \frac{\varepsilon}{2} \).
        \item
            Пусть \( \{ L_m \}_{m = 1}^{\infty} \) --- семейство множеств таких, что \( \mu(L_m) = 0 \).

            Тогда \( \mu \left( \bigcup\limits_{m = 1}^{\infty} L_m \right) = 0 \)
    \end{enumerate}
\end{example}

\begin{exercise}{}{}
    Доказать, что отрезок \( [a, b] : a \neq b \) не является
    множеством нулевой меры Лебега.
\end{exercise}

\begin{exercise}{}{}
    Привести пример континуального множества нулевой меры Лебега.
\end{exercise}

\begin{prop}{}{}
    Если \( \forall m \in \NN \mu(L_m) = 0 \),
    то \( \mu \left( \bigcup\limits_{m = 1}^{\infty} L_m \right) = 0 \).
\end{prop}

\begin{prop}{}{}
    Если \( L_1 \subseteq L \) и \( \mu(L) = 0 \), то \( \mu(L_1) = 0 \).
\end{prop}

\begin{prop}{}{}
    В определении множества \( L \) лебеговой меры ноль
    покрытие интервалами можно заменить на покрытие отрезками.
\end{prop}

\begin{thrm}{Критерий Лебега}{}
    \begin{gather*}
        f \in R[a, b]
        \\
        \Updownarrow
        \\
        \text{\( f \) ограничена и множество точек разрыва \( f \) имеет лебегову меру \( 0 \)}
    \end{gather*}
\end{thrm}
